\section*{Capítulo \ref{ch:b2:climate-proxies} --- Arquivos e indicadores climáticos}

\begin{solution}{exo:b2:climate-proxies:1}
(a) $\delta = (27.5 \times (-30) + 1350 \times 0)/(1350 + 27.5) = -825/1377.5 = -0.60$ por mil.
(b) O oceano conteria então toda a água: $-0.60$ por mil, \num{0.60} por mil mais leve que
hoje. (c) $27.5/1377.5 = \qty{2.0}{\percent}$.
\end{solution}

\begin{solution}{exo:b2:climate-proxies:2}
(a) $T = 16.1 - 4.76\,(2.50 - (0 - 0.27)) = 16.1 - 4.76 \times 2.77 = 16.1 - 13.19 =
\qty{2.9}{\celsius}$. (b) $T = 16.1 - 4.76\,(2.50 - 0.73) = 16.1 - 8.43 = \qty{7.7}{\celsius}$: a
mesma carapaça dá uma temperatura quase cinco graus mais alta se a água era mais pesada, e é por isso que
a composição da água deve ser conhecida. (c) A carapaça é $2.71 - 2.50 = 0.21$ por mil mais pesada,
logo sua água estava $4.76 \times 0.21 = \qty{1.0}{\celsius}$ mais fria, a \qty{1.9}{\celsius}.
\end{solution}

\begin{solution}{exo:b2:climate-proxies:3}
(a) $0.4\,\mathrm{e}^{0.9} = \qty{0.98}{mmol/mol}$ e $0.4\,\mathrm{e}^{1.8} =
\qty{2.42}{mmol/mol}$. (b) O fator é $\mathrm{e}^{0.9} = 2.46$; por grau $\mathrm{e}^{0.09}
= 1.094$, um aumento de \qty{9.4}{\percent} por grau. (c) $T = \ln(\text{Mg/Ca}/0.4)/0.09$:
$\ln 5/0.09 = \qty{17.9}{\celsius}$ e $\ln 2.5/0.09 = \qty{10.2}{\celsius}$.
\end{solution}

\begin{solution}{exo:b2:climate-proxies:4}
(a) $2668 - 2213 + 1 = 456$ anos. (b) Uma árvore isolada pode não ter um anel localmente, ou mostrar um anel
que não se consegue ler; comparar os padrões de anéis largos e estreitos em muitas árvores revela os
anéis em falta e fixa o ano de cada anel, e a concordância entre as árvores mede a
confiabilidade da datação. (c) No limite das árvores um só fator, o frio da estação de crescimento,
limita todas as árvores, de modo que elas partilham o mesmo padrão de anos bons e maus; numa baixada amena e
úmida nenhum fator domina, os padrões das árvores diferem, e a \omterm{def:b2:climate-proxies:dendro}{datação cruzada} pode falhar.
\end{solution}

\begin{solution}{exo:b2:climate-proxies:5}
(a) $125 \times 40/(3682.2 - 125) = 5000/3557.2 = 1.41$ por mil; $134 \times 40/3548.2 =
5360/3548.2 = 1.51$ por mil. (b) Com $H - h = 3552.2$: $3900/3552.2 = 1.10$, $4550/3552.2 = 1.28$
e $6500/3552.2 = 1.83$ por mil. (c) $\delta_{\mathrm{ice}} = -1.0 \times 3552.2/130 = -27.3$ por
mil, mais pesado que qualquer manto de gelo atual (Groenlândia $-30$ a $-35$, o manto de gelo norte-americano
$-40$, o gelo antártico chegando a $-50$). O balanço de massa com um gelo realista e a medição nas águas
intersticiais não concordam; a diferença de cerca de meio por mil é uma incerteza real do método.
\end{solution}

\begin{solution}{exo:b2:climate-proxies:6}
(a) $0.372\,\mathrm{e}^{0.0936T} = 0.4\,\mathrm{e}^{0.09T}$ dá $\mathrm{e}^{0.0036T} =
0.4/0.372 = 1.0753$, logo $T = \ln 1.0753/0.0036 = \qty{20.2}{\celsius}$. (b) A
\qty{2}{\celsius}: $0.4\,\mathrm{e}^{0.18} = 0.479$ e $0.372\,\mathrm{e}^{0.1872} = 0.449$, de modo que a
razão ajustada é $0.449/0.479 = 0.937$, \qty{6.3}{\percent} menor. (c) $T = \ln(0.449/0.4)/0.09 =
\ln 1.1225/0.09 = \qty{1.3}{\celsius}$: a lei prevista lê a carapaça cerca de
\qty{0.7}{\celsius} fria demais. Em água fria a escolha da calibração conta ao nível de um
grau.
\end{solution}

\begin{solution}{exo:b2:climate-proxies:7}
(a) Do Holoceno a GI-1, \qty{112.19}{m} de gelo para 2989 anos: $11\,219/2989 =
\qty{3.75}{cm}$ por ano; de GI-16 a GI-17, \qty{17.89}{m} para 1160 anos: $1789/1160 =
\qty{1.54}{cm}$ por ano. (b) O gelo
flui: à medida que cada camada é soterrada, ela se espalha lateralmente e se adelgaça, e as camadas mais profundas e mais antigas
foram adelgaçadas por mais tempo. (c) Quando as camadas ficam finas demais, um centímetro ou menos, os instrumentos
já não conseguem separar um ano do seguinte, e o \omterm{def:b2:climate-proxies:age-model}{modelo de idade} tem de usar pontos de amarração ou modelos de fluxo
em seu lugar.
\end{solution}

\begin{solution}{exo:b2:climate-proxies:8}
(a) A fração do intervalo é $(1550 - 1492.45)/(1604.64 - 1492.45) = 57.55/112.19 = 0.513$,
logo a idade é $11\,703 + 0.513 \times 2989 = 11\,703 + 1533 = \num{13236}$ anos b2k. (b) 1950 fica 50
anos antes de 2000: \num{13186} anos antes de 1950. (c) $50 + 0.513 \times (93 - 50) = 72$ anos.
\end{solution}

\begin{solution}{exo:b2:climate-proxies:9}
(a) Cada camada incerta contribui com meio ano: $2 \times 93 = 186$ e $2 \times 1287 = 2574$
camadas incertas. (b) $\frac12\sqrt{186} = 6.8$ anos e $\frac12\sqrt{2574} = 25.4$ anos.
(c) $93/6.8 = 14$ e $1287/25.4 = 51$. A razão entre $\frac12 n$ e $\frac12\sqrt{n}$ é
$\sqrt{n}$: ela cresce como a raiz quadrada do número de camadas incertas, de modo que a escolha entre
as duas regras pesa cada vez mais com a profundidade.
\end{solution}

\begin{solution}{exo:b2:climate-proxies:10}
(a) Na \omterm{prop:b2:isotope-geochemistry:mwl}{reta da água meteórica} $\delta$D varia oito vezes mais que $\delta^{18}$O, de modo que a
inclinação é $8 \times 0.8 = 6.4$ por mil de $\delta$D por grau. A queda é $441.7 - 396.6 = 45.1$
por mil, daí um resfriamento de $45.1/6.4 = \qty{7.0}{\celsius}$. (b) O oceano glacial era cerca de
$8 \times 1.0 = 8$ por mil mais pesado em $\delta$D; a neve formada a partir do oceano atual teria sido
mais leve na mesma quantidade, $-441.7 - 8 = -449.7$ por mil. A queda passa a $53.1$ por mil e o resfriamento a
$53.1/6.4 = \qty{8.3}{\celsius}$. (c) $9.39 - 0.09 = \qty{9.30}{\celsius}$, logo $45.1/9.30 = 4.85$
por mil por grau. A série publicada também é corrigida da altitude variável do manto
de gelo, e sua relação com a temperatura vem de modelos que simulam os isótopos da água, não apenas
da inclinação espacial.
\end{solution}

\begin{solution}{exo:b2:climate-proxies:11}
(a) Os logaritmos das razões são $-0.528$, $-0.223$, $0.191$, $0.501$, $0.904$, $1.227$; sua
média é $0.345$ e a temperatura média 14. Com $S_{TT} = \sum (T - 14)^2 = 280$ e $S_{T\ell} =
\sum (T - 14)(\ell - 0.345) = 24.93$, $A = 24.93/280 = \qty{0.0890}{\per\celsius}$ e $\ln B =
0.345 - 0.0890 \times 14 = -0.901$, logo $B = \qty{0.406}{mmol/mol}$. (b) Os resíduos são
$+0.017$, $-0.034$, $+0.023$, $-0.023$, $+0.025$ e $-0.009$, alguns por cento da razão, sem
tendência; o ajuste é próximo de $0.4\,\mathrm{e}^{0.09T}$. (c) $T = (\ln 1.30 + 0.901)/0.0890 = (0.262 +
0.901)/0.0890 = \qty{13.1}{\celsius}$. (d) $T = (\ln 0.45 + 0.901)/0.0890 = (-0.799 +
0.901)/0.0890 = \qty{1.2}{\celsius}$. Esse valor fica abaixo da amostra mais fria da calibração,
\qty{4}{\celsius}: a lei é extrapolada, e a baixas temperaturas sabe-se que a razão Mg/Ca das carapaças
se afasta da lei exponencial sob a influência do íon carbonato.
\end{solution}

\begin{solution}{exo:b2:climate-proxies:12}
(a) Por unidade de área do oceano, a água do oceano atual, $H$ a 0, reparte-se entre os reservatórios
$h_i$ a $\delta_i$ e o oceano restante $H - h$, com $h = \sum h_i$: $0 = \sum h_i\delta_i +
(H - h)\,\Delta\delta$. Depois da primeira etapa $\Delta\delta = 60 \times 30/(3682.2 - 60) =
1800/3622.2 = 0.50$ por mil; depois da segunda $(1800 + 70 \times 45)/(3682.2 - 130) =
4950/3552.2 = 1.39$ por mil. (b) $-4950/130 = -38.1$ por mil. (c) A regra é $1.39/130 = 0.0107$
por mil por metro, e dá $0.50/0.0107 = \qty{47}{m}$ para a primeira etapa em vez de
\qty{60}{m}: o primeiro gelo era mais pesado que a média, de modo que cada um de seus metros deslocou o oceano
menos. Uma regra linear calibrada num estado lê mal os outros quando a composição média do
gelo muda.
\end{solution}

\begin{solution}{pb:b2:climate-proxies:1}
\textbf{1.} $361.84 \times 10^6\ \unit{km^2} \times \qty{0.130}{km} = \qty{4.70e7}{km^3}$, cerca de 47
milhões de quilômetros cúbicos.
\textbf{2.} $130 \times 40/(3682.2 - 130) = 5200/3552.2 = 1.46$ por mil.
\textbf{3.} $5000/3557.2 = 1.41$ e $5360/3548.2 = 1.51$ por mil.
\textbf{4.} $4550/3552.2 = 1.28$ e $6500/3552.2 = 1.83$ por mil. A composição do gelo
desloca a resposta em \num{0.55} por mil, o nível do mar em \num{0.10}: o gelo pesa mais.
\textbf{5.} $\delta_{\mathrm{ice}} = -1.0 \times 3552.2/130 = -27.3$ por mil.
\textbf{6.} $40/3682.2 = 0.0109$ por mil por metro, que se arredonda para o \num{0.011} dos níveis do mar
datados.
\textbf{7.} Com $\delta^{18}\mathrm{O}_{\mathrm{sw}}$ constante, $\mathrm{d}T = -4.76\,\mathrm{d}\delta^{18}
\mathrm{O}_{\mathrm{b}}$, logo $\mathrm{d}T = -1$ dá $\mathrm{d}\delta^{18}\mathrm{O}_{\mathrm{b}} =
1/4.76 = 0.21$ por mil.
\textbf{8.} $\delta^{18}\mathrm{O}_{\mathrm{b}} = (16.1 - T)/4.76 + \delta^{18}\mathrm{O}_{\mathrm{sw}}
- 0.27 = 14.1/4.76 - 0.27 = 2.96 - 0.27 = 2.69$ por mil.
\textbf{9.} $17.1/4.76 + 1.0 - 0.27 = 3.59 + 0.73 = 4.32$ por mil. A variação é $4.32 - 2.69 =
1.63$ por mil: \num{1.00} vindo da água e $3.0 \times 0.21 = 0.63$ do resfriamento de
\qty{3}{\celsius}.
\textbf{10.} $T = \ln(2.0/0.4)/0.09 = \ln 5/0.09 = \qty{17.9}{\celsius}$.
\textbf{11.} $0.18/0.1 = \qty{1.8}{\celsius}$ de resfriamento.
\textbf{12.} A razão Mg/Ca é uma razão entre dois elementos na calcita, fixada pela temperatura da
calcificação; não envolve os isótopos do oxigênio, e por isso ignora a composição da
água que o gelo modifica.
\textbf{13.} $4.92 - 3.23 = 1.69$ por mil.
\textbf{14.} Resolvendo $\Delta = h \times 40/(H - h)$ em $h$: $h = H\Delta/(\Delta + 40) = 3682.2
\times 1.69/41.69 = \qty{149}{m}$, mais que os 125 a \qty{134}{m} dos indícios do nível do mar:
parte da variação tem de ser resfriamento.
\textbf{15.} $1.69 - 1.46 = 0.23$ por mil, e $4.76 \times 0.23 = \qty{1.1}{\celsius}$.
\textbf{16.} $1.69 - 1.00 = 0.69$ por mil, e $4.76 \times 0.69 = \qty{3.3}{\celsius}$ (3.28).
\textbf{17.} $0.23/0.25 = \qty{0.9}{\celsius}$ e $0.69/0.25 = \qty{2.8}{\celsius}$.
\textbf{18.} A Água Profunda do Atlântico Norte resfriou de 2--\qty{4}{\celsius} para \qty{-1.1}{\celsius},
ou seja 3.1 a \qty{5.1}{\celsius}; a Água de Fundo Antártica de 0--\qty{1}{\celsius} para
\qty{-1.9}{\celsius}, ou seja 1.9 a \qty{2.9}{\celsius}; a razão Mg/Ca do Atlântico Noroeste dá
\qty{1.8}{\celsius}. Todos ficam entre cerca de 2 e \qty{5}{\celsius}, em torno dos 2.8 a
\qty{3.3}{\celsius} da correção pelas águas intersticiais, e acima do \qty{1}{\celsius} do balanço
de massa com gelo a $-40$ por mil.
\textbf{19.} $-396.6 - (-441.7) = 45.1$ por mil.
\textbf{20.} $8 \times 0.8 = 6.4$ por mil por grau, logo $45.1/6.4 = \qty{7.0}{\celsius}$.
\textbf{21.} $9.39 - 0.09 = \qty{9.30}{\celsius}$, e $45.1/9.30 = 4.85$ por mil por grau.
\textbf{22.} Dome C, \qty{9.3}{\celsius}, mais que a superfície global, 5 a \qty{7}{\celsius}, mais
que o oceano profundo, cerca de \qty{3}{\celsius}.
\textbf{23.} A água profunda já é fria hoje, 0 a \qty{4}{\celsius}, e a água do mar não pode
resfriar muito abaixo de seu ponto de congelamento, que a água de fundo glacial quase atingiu, a
\qty{-1.9}{\celsius}: o oceano profundo só tinha alguns graus a perder, enquanto o ar sobre a Antártida
não tem esse piso.
\textbf{24.} Com a correção pelas águas intersticiais, o oceano profundo estava cerca de \textbf{\qty{3.3}{\celsius}
mais frio} no Último Máximo Glacial do que hoje, perto dos 2 a \qty{5}{\celsius} das estimativas
independentes da água profunda.
\end{solution}
